\chapter{Geometry and Dynamics of Learning}

\section{Loss Landscapes}

Training a neural network can be viewed as navigating a high-dimensional loss landscape.

\medskip

\noindent
\textbf{Loss function.}  
Given parameters $\theta \in \mathbb{R}^d$, the loss defines a function:
\[
L(\theta): \mathbb{R}^d \to \mathbb{R}.
\]

\medskip

\noindent
\textbf{Landscape interpretation.}
\begin{itemize}
    \item Points correspond to parameter configurations
    \item Height corresponds to loss value
\end{itemize}

\medskip

\noindent
\textbf{Critical points.}
\begin{itemize}
    \item Local minima: $\nabla L(\theta) = 0$, positive curvature
    \item Saddle points: mixed curvature directions
\end{itemize}

\medskip

\noindent
\textbf{High-dimensional geometry.}
\begin{itemize}
    \item Saddle points are more prevalent than poor local minima
    \item Many directions are nearly flat
\end{itemize}

\medskip

\noindent
\textbf{Flat vs sharp minima.}
\begin{itemize}
    \item Flat minima: robust to perturbations
    \item Sharp minima: sensitive to parameter changes
\end{itemize}

\medskip

\noindent
\textbf{Visualization (conceptual).}  
Low-dimensional projections reveal valleys, ridges, and plateaus.

\medskip

\noindent
\textbf{Insight.}  
The geometry of the loss landscape strongly influences optimization and generalization.

\section{Neural Tangent Kernel (Intuition)}

The neural tangent kernel (NTK) provides a linearized view of neural network training.

\medskip

\noindent
\textbf{Linearization.}  
Around initialization $\theta_0$, the network can be approximated as:
\[
f_\theta(x) \approx f_{\theta_0}(x) + \nabla_\theta f_{\theta_0}(x)^\top (\theta - \theta_0).
\]

\medskip

\noindent
\textbf{Neural tangent kernel.}
\[
K(x, x') = \nabla_\theta f_{\theta_0}(x)^\top \nabla_\theta f_{\theta_0}(x').
\]

\medskip

\noindent
\textbf{Interpretation.}
\begin{itemize}
    \item Defines a kernel induced by the network at initialization
    \item Training resembles kernel regression with kernel $K$
\end{itemize}

\medskip

\noindent
\textbf{Infinite-width limit.}
\begin{itemize}
    \item NTK becomes deterministic
    \item Training dynamics become linear
\end{itemize}

\medskip

\noindent
\textbf{Insight.}  
The NTK connects deep learning with classical kernel methods.

\section{Continuous-Time View of Training}

Gradient-based optimization can be viewed as a continuous-time dynamical system.

\medskip

\noindent
\textbf{Gradient descent update.}
\[
\theta_{t+1} = \theta_t - \eta \nabla L(\theta_t).
\]

\medskip

\noindent
\textbf{Continuous-time limit.}  
As $\eta \to 0$, we obtain a differential equation:
\[
\frac{d\theta(t)}{dt} = -\nabla L(\theta(t)).
\]

\medskip

\noindent
\textbf{Gradient flow.}
\begin{itemize}
    \item Parameters evolve along the steepest descent direction
    \item Trajectories follow curves in parameter space
\end{itemize}

\medskip

\noindent
\textbf{Energy interpretation.}  
The loss decreases along trajectories:
\[
\frac{d}{dt} L(\theta(t)) = -\|\nabla L(\theta(t))\|^2 \leq 0.
\]

\medskip

\noindent
\textbf{Extension to stochastic dynamics.}  
Stochastic gradient descent introduces noise:
\[
d\theta(t) = -\nabla L(\theta(t)) dt + \Sigma^{1/2} dW_t,
\]
where $W_t$ is Brownian motion.

\medskip

\noindent
\textbf{Insight.}  
Training can be understood as a continuous flow shaped by both gradients and noise.

\section{Dynamical Systems Perspective}

Neural network training and representation learning can be viewed through the lens of dynamical systems.

\medskip

\noindent
\textbf{Parameter dynamics.}
\begin{itemize}
    \item Evolution of $\theta(t)$ under optimization
\end{itemize}

\medskip

\noindent
\textbf{Representation dynamics.}
\begin{itemize}
    \item Evolution of features $f_{\theta(t)}(x)$ during training
\end{itemize}

\medskip

\noindent
\textbf{Stability.}
\begin{itemize}
    \item Stable fixed points correspond to minima
    \item Stability depends on curvature and noise
\end{itemize}

\medskip

\noindent
\textbf{Attractors.}
\begin{itemize}
    \item Optimization may converge to particular regions of parameter space
\end{itemize}

\medskip

\noindent
\textbf{Connection to residual networks.}  
Deep networks resemble discretized dynamical systems.

\medskip

\noindent
\textbf{Feature evolution.}
\begin{itemize}
    \item Early training: rapid feature changes
    \item Later stages: fine adjustments
\end{itemize}

\medskip

\noindent
\textbf{Insight.}  
Learning is a dynamic process that evolves both parameters and representations over time.

\section{A Unifying Perspective}

The geometry and dynamics of learning provide a unified framework:

\begin{itemize}
    \item \textbf{Geometry:} structure of the loss landscape
    \item \textbf{Dynamics:} trajectories of optimization
    \item \textbf{Linearization:} NTK approximation
    \item \textbf{Stochasticity:} role of noise in training
\end{itemize}

\medskip

\noindent
\textbf{Core idea.}  
Understanding learning requires analyzing both the structure of the objective and the dynamics of optimization.

\medskip

\noindent
\textbf{Key insight.}  
Modern deep learning lies at the intersection of geometry, probability, and dynamical systems.

\section{Outlook}

This chapter examined the geometry and dynamics of learning:
\begin{itemize}
    \item loss landscapes and their structure,
    \item neural tangent kernel perspective,
    \item continuous-time view of optimization,
    \item dynamical systems interpretation.
\end{itemize}

\medskip

These perspectives provide a deeper theoretical understanding of how learning occurs in modern machine learning systems.

\medskip

\noindent
\textbf{Guiding principle.}  
Learning is both a geometric problem and a dynamical process, shaped by the interplay between model structure and optimization.